% =====================================================================
% Chapter 1: Bloch's theorem and band folding
% Theory session, dispersion branch (Nov 2026)
% Requires: amsmath, amssymb
% Include in the main file with  % =====================================================================
% Chapter 1: Bloch's theorem and band folding
% Theory session, dispersion branch (Nov 2026)
% Requires: amsmath, amssymb
% Include in the main file with  % =====================================================================
% Chapter 1: Bloch's theorem and band folding
% Theory session, dispersion branch (Nov 2026)
% Requires: amsmath, amssymb
% Include in the main file with  % =====================================================================
% Chapter 1: Bloch's theorem and band folding
% Theory session, dispersion branch (Nov 2026)
% Requires: amsmath, amssymb
% Include in the main file with  \include{chapter_01_bloch_theorem}
% =====================================================================

\chapter{Bloch's theorem and band folding}
\label{ch:bloch}

Numbers in this chapter refer to the production configuration of the Bloch
unit cell: steel plate, thickness $h = 5\,\mathrm{mm}$, square unit cell of
size $a = 0.1\,\mathrm{m}$, Reissner--Mindlin (RM) plate theory, bands below
$5\,\mathrm{kHz}$. Plate data: $\rho h \approx 39.3\,\mathrm{kg/m^2}$,
$D \approx 2.3\,\mathrm{kN\,m}$, plate mass per cell
$\rho h a^2 \approx 0.39\,\mathrm{kg}$.

% ---------------------------------------------------------------------
\section{What Bloch's theorem says}

A periodic structure repeats along the lattice vectors
$\mathbf{a}_1 = (L_x, 0)$ and $\mathbf{a}_2 = (0, L_y)$. A lattice vector is
any $\mathbf{R} = m\,\mathbf{a}_1 + n\,\mathbf{a}_2$ with integers $m, n$.
Bloch's theorem states that every free wave can be written in the form
\begin{equation}
  u(\mathbf{x} + \mathbf{R}) = e^{i\,\mathbf{k}\cdot\mathbf{R}}\, u(\mathbf{x}),
  \label{eq:bloch}
\end{equation}
or equivalently $u(\mathbf{x}) = e^{i\mathbf{k}\cdot\mathbf{x}}\,p(\mathbf{x})$
with $p$ periodic in the cell. Moving one cell over does not change the shape
of the wave, it only multiplies it by a phase. Therefore one cell, with the
conditions ``right edge = left edge $\cdot\, e^{i k_x L_x}$'' and
``top edge = bottom edge $\cdot\, e^{i k_y L_y}$'', describes the whole
infinite structure.

% ---------------------------------------------------------------------
\section{Why it holds: the argument in four steps}

\paragraph{Step 1: the structure is unchanged by a lattice shift.}
Define the translation operator $\mathcal{T}_{\mathbf{R}}$ by
$(\mathcal{T}_{\mathbf{R}} u)(\mathbf{x}) = u(\mathbf{x} + \mathbf{R})$.
The equations of motion of the plate have periodic coefficients (the resonator
sits at the same place in every cell). It therefore makes no difference
whether we first shift and then apply the equations, or the other way round:
the two operators \emph{commute}.

\paragraph{Step 2: a shifted solution is again a solution.}
If $u$ is a free wave at frequency $\omega$, then $\mathcal{T}_{\mathbf{R}}u$
is also a free wave at the same $\omega$. Operators that commute share
eigenfunctions. Within the space of solutions at that $\omega$ we can
therefore \emph{choose} solutions that $\mathcal{T}_{\mathbf{R}}$ only
multiplies by a number:
\begin{equation}
  \mathcal{T}_{\mathbf{R}} u = c(\mathbf{R})\, u .
\end{equation}

\paragraph{Step 3: the multiplier must be an exponential.}
Shifting by $\mathbf{R}_1$ and then by $\mathbf{R}_2$ equals shifting by
$\mathbf{R}_1 + \mathbf{R}_2$, so
$c(\mathbf{R}_1 + \mathbf{R}_2) = c(\mathbf{R}_1)\,c(\mathbf{R}_2)$.
The only continuous functions with this property are exponentials:
\begin{equation}
  c(\mathbf{R}) = e^{i\,\mathbf{k}\cdot\mathbf{R}} .
\end{equation}

\paragraph{Step 4: what $\mathbf{k}$ can be.}
\begin{itemize}
  \item If the wave stays bounded over the infinite lattice, then $|c| = 1$
        and $\mathbf{k}$ is \emph{real}. These are the propagating waves,
        and they form the bands.
  \item If $|c| \neq 1$, then $\mathbf{k}$ is \emph{complex}: the wave grows
        or decays from cell to cell (evanescent wave).
\end{itemize}
Keep the second case in mind: inside a band gap no real-$\mathbf{k}$
solutions exist, but complex-$\mathbf{k}$ solutions do. The proof itself does
not restrict $\mathbf{k}$ to real values; restricting to real $\mathbf{k}$ is
our choice, because we want to know which waves propagate.

% ---------------------------------------------------------------------
\section{Consequences}

\begin{itemize}
  \item \textbf{One cell is enough, but only for an infinite lattice.} A finite
        plate has edges; a shift no longer leaves it unchanged, and Step 1
        fails. The band diagram describes the \emph{material}, not the finite
        $1.1 \times 1.0\,\mathrm{m}$ plate.
  \item \textbf{$\mathbf{k}$ only matters modulo $\mathbf{G}$.} The phase
        $e^{i\mathbf{k}\cdot\mathbf{R}}$ does not change when we add a
        reciprocal lattice vector
        $\mathbf{G} = 2\pi\,(m/L_x,\; n/L_y)$, because
        $\mathbf{G}\cdot\mathbf{R}$ is a multiple of $2\pi$. This is the origin
        of the Brillouin zone and of folding.
  \item \textbf{At each $\mathbf{k}$ the problem is finite.} The cell has
        finitely many degrees of freedom, so it has a discrete list of
        frequencies. Sorted, they are bands $1, 2, 3, \dots$ at that
        $\mathbf{k}$.
\end{itemize}

\paragraph{Brillouin zone and symmetry points.}
All distinct wavevectors lie in the first Brillouin zone
$-\pi/L_x \le k_x \le \pi/L_x$, $-\pi/L_y \le k_y \le \pi/L_y$. Reciprocity
($\omega(\mathbf{k}) = \omega(-\mathbf{k})$) and the mirror symmetries of the
cell reduce this to the irreducible Brillouin zone, and band diagrams follow its
edge.

\begin{center}
\begin{tabular}{|l|l|l|l|}
\hline
Point & $\mathbf{k}$ & Phase per cell & Picture \\
\hline
$\Gamma$ & $(0,0)$ & $1,\ 1$ & all cells move identically \\
$X$ & $(\pi/L_x, 0)$ & $-1$ in $x$ & neighbours in $x$ move oppositely \\
$M$ & $(\pi/L_x, \pi/L_y)$ & $-1,\ -1$ & checkerboard \\
$Y$ & $(0, \pi/L_y)$ & $-1$ in $y$ & only needed for rectangular cells \\
\hline
\end{tabular}
\end{center}
Path: $\Gamma$--$X$--$M$--$\Gamma$ for a square cell,
$\Gamma$--$X$--$M$--$Y$--$\Gamma$ for a rectangular cell.

% ---------------------------------------------------------------------
\section{Band folding with our numbers}

For the uniform plate (no cells) each plane wave $e^{i\boldsymbol{\kappa}\cdot\mathbf{x}}$
has one frequency, set by the physical wavenumber $\kappa = |\boldsymbol{\kappa}|$.
Thin-plate (Kirchhoff) estimate for the $5\,\mathrm{mm}$ plate:
\begin{equation}
  f \approx 1.215\,\kappa^2\ \mathrm{Hz}\qquad(\kappa \text{ in rad/m}).
\end{equation}
Described with cells of size $a$, the Bloch wavevector only sees
$\boldsymbol{\kappa}$ modulo $\mathbf{G}$. At a fixed $\mathbf{k}$ the cell
therefore carries all waves $\boldsymbol{\kappa} = \mathbf{k} + \mathbf{G}$,
each with its own frequency. Folding is a relabelling, nothing physical
changes.

\begin{center}
\begin{tabular}{|l|l|l|l|l|}
\hline
Point & $\kappa$ folding there & Kirchhoff & RM (FE) & Multiplicity \\
\hline
$X$ & $\pi/a = 31.4$ & $1199\,\mathrm{Hz}$ & $\approx 1190\,\mathrm{Hz}$ & 2: $(\pm\pi/a, 0)$ \\
$M$ & $\sqrt{2}\,\pi/a = 44.4$ & $2398\,\mathrm{Hz}$ & $\approx 2380\,\mathrm{Hz}$ & 4: $(\pm\pi/a, \pm\pi/a)$ \\
$\Gamma$ & $2\pi/a = 62.8$ & $4797\,\mathrm{Hz}$ & $\approx 4730\,\mathrm{Hz}$ & 4: $(\pm 2\pi/a, 0), (0, \pm 2\pi/a)$ \\
\hline
\end{tabular}
\end{center}
The difference between Kirchhoff and RM grows with frequency
($0.8\,\%$ at $X$, $1.4\,\%$ at $\Gamma$), because shear deformation and
rotary inertia grow with $\kappa h$. This is also why the Kirchhoff
plane-wave model was $1.6\,\%$ off in the stage~4 validation.

% ---------------------------------------------------------------------
\section{Check questions with answers}

\paragraph{Question 1.} \emph{In Step 2, why can we only say that we can
choose solutions that $\mathcal{T}_{\mathbf{R}}$ multiplies by a number,
instead of claiming that every solution is like that?}

\emph{Answer.} Bloch waves are a basis we choose; not every solution is one.
Example: at a generic frequency the plate carries $e^{+i\kappa x}$ and
$e^{-i\kappa x}$, which have \emph{different} phases per cell,
$e^{\pm i\kappa a}$. Their sum $\cos(\kappa x)$ solves the equations at the
same $\omega$, but shifting it by $a$ does not give a multiple of
$\cos(\kappa x)$. It is a valid solution, but not a Bloch wave.

\emph{At $X$ this changes:} $e^{+i\pi x/a}$ and $e^{-i\pi x/a}$ both have the
phase $e^{\pm i\pi} = -1$ per cell. They differ by exactly one $\mathbf{G}$
and are \emph{the same} Bloch $\mathbf{k}$. Every combination of them is
therefore also a Bloch wave at $X$, including $\cos(\pi x/a)$ and
$\sin(\pi x/a)$. This is why the standing waves A and B used in
Chapter~\ref{ch:two_level} are legitimate Bloch modes, and why a resonator can
split them.

\paragraph{Question 2.} \emph{At $\Gamma$ the bare cell has a band at about
$4.73\,\mathrm{kHz}$, although ``$\Gamma$ means infinite wavelength''. What is
that wave physically?}

\emph{Answer.} The Bloch wavevector is $\mathbf{k} = 0$, but the physical
wavenumber is $\kappa = |\mathbf{k} + \mathbf{G}| = 2\pi/a$. Both describe the
same mode because $\mathbf{k}$ only matters modulo $\mathbf{G}$. Physically
it is a wave with wavelength $\lambda_\mathrm{wave} = a = 0.1\,\mathrm{m}$:
exactly one full wiggle per cell, phase per cell $e^{i2\pi} = 1$. All cells
look identical, which is why it sits at $\Gamma$. It appears four times, once
for each direction $\pm x$, $\pm y$.

\paragraph{Question 3.} \emph{What breaks first if Bloch's theorem is applied
to the finite $1.1 \times 1.0\,\mathrm{m}$ plate, and what goes wrong
physically?}

\emph{Answer.} Step 1 already fails: a shifted solution no longer satisfies
the edge conditions, so the operators no longer commute. Physically, waves
reflect at the free edges, travelling waves become standing waves, and the
finite plate has discrete modes. Two consequences (treated later):
\begin{itemize}
  \item the modes of the finite plate roughly \emph{sample} the bands at
        discrete wavenumbers, about $n\pi/L$;
  \item edges can carry localized modes, sometimes even inside a band gap.
\end{itemize}

% =====================================================================

\chapter{Bloch's theorem and band folding}
\label{ch:bloch}

Numbers in this chapter refer to the production configuration of the Bloch
unit cell: steel plate, thickness $h = 5\,\mathrm{mm}$, square unit cell of
size $a = 0.1\,\mathrm{m}$, Reissner--Mindlin (RM) plate theory, bands below
$5\,\mathrm{kHz}$. Plate data: $\rho h \approx 39.3\,\mathrm{kg/m^2}$,
$D \approx 2.3\,\mathrm{kN\,m}$, plate mass per cell
$\rho h a^2 \approx 0.39\,\mathrm{kg}$.

% ---------------------------------------------------------------------
\section{What Bloch's theorem says}

A periodic structure repeats along the lattice vectors
$\mathbf{a}_1 = (L_x, 0)$ and $\mathbf{a}_2 = (0, L_y)$. A lattice vector is
any $\mathbf{R} = m\,\mathbf{a}_1 + n\,\mathbf{a}_2$ with integers $m, n$.
Bloch's theorem states that every free wave can be written in the form
\begin{equation}
  u(\mathbf{x} + \mathbf{R}) = e^{i\,\mathbf{k}\cdot\mathbf{R}}\, u(\mathbf{x}),
  \label{eq:bloch}
\end{equation}
or equivalently $u(\mathbf{x}) = e^{i\mathbf{k}\cdot\mathbf{x}}\,p(\mathbf{x})$
with $p$ periodic in the cell. Moving one cell over does not change the shape
of the wave, it only multiplies it by a phase. Therefore one cell, with the
conditions ``right edge = left edge $\cdot\, e^{i k_x L_x}$'' and
``top edge = bottom edge $\cdot\, e^{i k_y L_y}$'', describes the whole
infinite structure.

% ---------------------------------------------------------------------
\section{Why it holds: the argument in four steps}

\paragraph{Step 1: the structure is unchanged by a lattice shift.}
Define the translation operator $\mathcal{T}_{\mathbf{R}}$ by
$(\mathcal{T}_{\mathbf{R}} u)(\mathbf{x}) = u(\mathbf{x} + \mathbf{R})$.
The equations of motion of the plate have periodic coefficients (the resonator
sits at the same place in every cell). It therefore makes no difference
whether we first shift and then apply the equations, or the other way round:
the two operators \emph{commute}.

\paragraph{Step 2: a shifted solution is again a solution.}
If $u$ is a free wave at frequency $\omega$, then $\mathcal{T}_{\mathbf{R}}u$
is also a free wave at the same $\omega$. Operators that commute share
eigenfunctions. Within the space of solutions at that $\omega$ we can
therefore \emph{choose} solutions that $\mathcal{T}_{\mathbf{R}}$ only
multiplies by a number:
\begin{equation}
  \mathcal{T}_{\mathbf{R}} u = c(\mathbf{R})\, u .
\end{equation}

\paragraph{Step 3: the multiplier must be an exponential.}
Shifting by $\mathbf{R}_1$ and then by $\mathbf{R}_2$ equals shifting by
$\mathbf{R}_1 + \mathbf{R}_2$, so
$c(\mathbf{R}_1 + \mathbf{R}_2) = c(\mathbf{R}_1)\,c(\mathbf{R}_2)$.
The only continuous functions with this property are exponentials:
\begin{equation}
  c(\mathbf{R}) = e^{i\,\mathbf{k}\cdot\mathbf{R}} .
\end{equation}

\paragraph{Step 4: what $\mathbf{k}$ can be.}
\begin{itemize}
  \item If the wave stays bounded over the infinite lattice, then $|c| = 1$
        and $\mathbf{k}$ is \emph{real}. These are the propagating waves,
        and they form the bands.
  \item If $|c| \neq 1$, then $\mathbf{k}$ is \emph{complex}: the wave grows
        or decays from cell to cell (evanescent wave).
\end{itemize}
Keep the second case in mind: inside a band gap no real-$\mathbf{k}$
solutions exist, but complex-$\mathbf{k}$ solutions do. The proof itself does
not restrict $\mathbf{k}$ to real values; restricting to real $\mathbf{k}$ is
our choice, because we want to know which waves propagate.

% ---------------------------------------------------------------------
\section{Consequences}

\begin{itemize}
  \item \textbf{One cell is enough, but only for an infinite lattice.} A finite
        plate has edges; a shift no longer leaves it unchanged, and Step 1
        fails. The band diagram describes the \emph{material}, not the finite
        $1.1 \times 1.0\,\mathrm{m}$ plate.
  \item \textbf{$\mathbf{k}$ only matters modulo $\mathbf{G}$.} The phase
        $e^{i\mathbf{k}\cdot\mathbf{R}}$ does not change when we add a
        reciprocal lattice vector
        $\mathbf{G} = 2\pi\,(m/L_x,\; n/L_y)$, because
        $\mathbf{G}\cdot\mathbf{R}$ is a multiple of $2\pi$. This is the origin
        of the Brillouin zone and of folding.
  \item \textbf{At each $\mathbf{k}$ the problem is finite.} The cell has
        finitely many degrees of freedom, so it has a discrete list of
        frequencies. Sorted, they are bands $1, 2, 3, \dots$ at that
        $\mathbf{k}$.
\end{itemize}

\paragraph{Brillouin zone and symmetry points.}
All distinct wavevectors lie in the first Brillouin zone
$-\pi/L_x \le k_x \le \pi/L_x$, $-\pi/L_y \le k_y \le \pi/L_y$. Reciprocity
($\omega(\mathbf{k}) = \omega(-\mathbf{k})$) and the mirror symmetries of the
cell reduce this to the irreducible Brillouin zone, and band diagrams follow its
edge.

\begin{center}
\begin{tabular}{|l|l|l|l|}
\hline
Point & $\mathbf{k}$ & Phase per cell & Picture \\
\hline
$\Gamma$ & $(0,0)$ & $1,\ 1$ & all cells move identically \\
$X$ & $(\pi/L_x, 0)$ & $-1$ in $x$ & neighbours in $x$ move oppositely \\
$M$ & $(\pi/L_x, \pi/L_y)$ & $-1,\ -1$ & checkerboard \\
$Y$ & $(0, \pi/L_y)$ & $-1$ in $y$ & only needed for rectangular cells \\
\hline
\end{tabular}
\end{center}
Path: $\Gamma$--$X$--$M$--$\Gamma$ for a square cell,
$\Gamma$--$X$--$M$--$Y$--$\Gamma$ for a rectangular cell.

% ---------------------------------------------------------------------
\section{Band folding with our numbers}

For the uniform plate (no cells) each plane wave $e^{i\boldsymbol{\kappa}\cdot\mathbf{x}}$
has one frequency, set by the physical wavenumber $\kappa = |\boldsymbol{\kappa}|$.
Thin-plate (Kirchhoff) estimate for the $5\,\mathrm{mm}$ plate:
\begin{equation}
  f \approx 1.215\,\kappa^2\ \mathrm{Hz}\qquad(\kappa \text{ in rad/m}).
\end{equation}
Described with cells of size $a$, the Bloch wavevector only sees
$\boldsymbol{\kappa}$ modulo $\mathbf{G}$. At a fixed $\mathbf{k}$ the cell
therefore carries all waves $\boldsymbol{\kappa} = \mathbf{k} + \mathbf{G}$,
each with its own frequency. Folding is a relabelling, nothing physical
changes.

\begin{center}
\begin{tabular}{|l|l|l|l|l|}
\hline
Point & $\kappa$ folding there & Kirchhoff & RM (FE) & Multiplicity \\
\hline
$X$ & $\pi/a = 31.4$ & $1199\,\mathrm{Hz}$ & $\approx 1190\,\mathrm{Hz}$ & 2: $(\pm\pi/a, 0)$ \\
$M$ & $\sqrt{2}\,\pi/a = 44.4$ & $2398\,\mathrm{Hz}$ & $\approx 2380\,\mathrm{Hz}$ & 4: $(\pm\pi/a, \pm\pi/a)$ \\
$\Gamma$ & $2\pi/a = 62.8$ & $4797\,\mathrm{Hz}$ & $\approx 4730\,\mathrm{Hz}$ & 4: $(\pm 2\pi/a, 0), (0, \pm 2\pi/a)$ \\
\hline
\end{tabular}
\end{center}
The difference between Kirchhoff and RM grows with frequency
($0.8\,\%$ at $X$, $1.4\,\%$ at $\Gamma$), because shear deformation and
rotary inertia grow with $\kappa h$. This is also why the Kirchhoff
plane-wave model was $1.6\,\%$ off in the stage~4 validation.

% ---------------------------------------------------------------------
\section{Check questions with answers}

\paragraph{Question 1.} \emph{In Step 2, why can we only say that we can
choose solutions that $\mathcal{T}_{\mathbf{R}}$ multiplies by a number,
instead of claiming that every solution is like that?}

\emph{Answer.} Bloch waves are a basis we choose; not every solution is one.
Example: at a generic frequency the plate carries $e^{+i\kappa x}$ and
$e^{-i\kappa x}$, which have \emph{different} phases per cell,
$e^{\pm i\kappa a}$. Their sum $\cos(\kappa x)$ solves the equations at the
same $\omega$, but shifting it by $a$ does not give a multiple of
$\cos(\kappa x)$. It is a valid solution, but not a Bloch wave.

\emph{At $X$ this changes:} $e^{+i\pi x/a}$ and $e^{-i\pi x/a}$ both have the
phase $e^{\pm i\pi} = -1$ per cell. They differ by exactly one $\mathbf{G}$
and are \emph{the same} Bloch $\mathbf{k}$. Every combination of them is
therefore also a Bloch wave at $X$, including $\cos(\pi x/a)$ and
$\sin(\pi x/a)$. This is why the standing waves A and B used in
Chapter~\ref{ch:two_level} are legitimate Bloch modes, and why a resonator can
split them.

\paragraph{Question 2.} \emph{At $\Gamma$ the bare cell has a band at about
$4.73\,\mathrm{kHz}$, although ``$\Gamma$ means infinite wavelength''. What is
that wave physically?}

\emph{Answer.} The Bloch wavevector is $\mathbf{k} = 0$, but the physical
wavenumber is $\kappa = |\mathbf{k} + \mathbf{G}| = 2\pi/a$. Both describe the
same mode because $\mathbf{k}$ only matters modulo $\mathbf{G}$. Physically
it is a wave with wavelength $\lambda_\mathrm{wave} = a = 0.1\,\mathrm{m}$:
exactly one full wiggle per cell, phase per cell $e^{i2\pi} = 1$. All cells
look identical, which is why it sits at $\Gamma$. It appears four times, once
for each direction $\pm x$, $\pm y$.

\paragraph{Question 3.} \emph{What breaks first if Bloch's theorem is applied
to the finite $1.1 \times 1.0\,\mathrm{m}$ plate, and what goes wrong
physically?}

\emph{Answer.} Step 1 already fails: a shifted solution no longer satisfies
the edge conditions, so the operators no longer commute. Physically, waves
reflect at the free edges, travelling waves become standing waves, and the
finite plate has discrete modes. Two consequences (treated later):
\begin{itemize}
  \item the modes of the finite plate roughly \emph{sample} the bands at
        discrete wavenumbers, about $n\pi/L$;
  \item edges can carry localized modes, sometimes even inside a band gap.
\end{itemize}

% =====================================================================

\chapter{Bloch's theorem and band folding}
\label{ch:bloch}

Numbers in this chapter refer to the production configuration of the Bloch
unit cell: steel plate, thickness $h = 5\,\mathrm{mm}$, square unit cell of
size $a = 0.1\,\mathrm{m}$, Reissner--Mindlin (RM) plate theory, bands below
$5\,\mathrm{kHz}$. Plate data: $\rho h \approx 39.3\,\mathrm{kg/m^2}$,
$D \approx 2.3\,\mathrm{kN\,m}$, plate mass per cell
$\rho h a^2 \approx 0.39\,\mathrm{kg}$.

% ---------------------------------------------------------------------
\section{What Bloch's theorem says}

A periodic structure repeats along the lattice vectors
$\mathbf{a}_1 = (L_x, 0)$ and $\mathbf{a}_2 = (0, L_y)$. A lattice vector is
any $\mathbf{R} = m\,\mathbf{a}_1 + n\,\mathbf{a}_2$ with integers $m, n$.
Bloch's theorem states that every free wave can be written in the form
\begin{equation}
  u(\mathbf{x} + \mathbf{R}) = e^{i\,\mathbf{k}\cdot\mathbf{R}}\, u(\mathbf{x}),
  \label{eq:bloch}
\end{equation}
or equivalently $u(\mathbf{x}) = e^{i\mathbf{k}\cdot\mathbf{x}}\,p(\mathbf{x})$
with $p$ periodic in the cell. Moving one cell over does not change the shape
of the wave, it only multiplies it by a phase. Therefore one cell, with the
conditions ``right edge = left edge $\cdot\, e^{i k_x L_x}$'' and
``top edge = bottom edge $\cdot\, e^{i k_y L_y}$'', describes the whole
infinite structure.

% ---------------------------------------------------------------------
\section{Why it holds: the argument in four steps}

\paragraph{Step 1: the structure is unchanged by a lattice shift.}
Define the translation operator $\mathcal{T}_{\mathbf{R}}$ by
$(\mathcal{T}_{\mathbf{R}} u)(\mathbf{x}) = u(\mathbf{x} + \mathbf{R})$.
The equations of motion of the plate have periodic coefficients (the resonator
sits at the same place in every cell). It therefore makes no difference
whether we first shift and then apply the equations, or the other way round:
the two operators \emph{commute}.

\paragraph{Step 2: a shifted solution is again a solution.}
If $u$ is a free wave at frequency $\omega$, then $\mathcal{T}_{\mathbf{R}}u$
is also a free wave at the same $\omega$. Operators that commute share
eigenfunctions. Within the space of solutions at that $\omega$ we can
therefore \emph{choose} solutions that $\mathcal{T}_{\mathbf{R}}$ only
multiplies by a number:
\begin{equation}
  \mathcal{T}_{\mathbf{R}} u = c(\mathbf{R})\, u .
\end{equation}

\paragraph{Step 3: the multiplier must be an exponential.}
Shifting by $\mathbf{R}_1$ and then by $\mathbf{R}_2$ equals shifting by
$\mathbf{R}_1 + \mathbf{R}_2$, so
$c(\mathbf{R}_1 + \mathbf{R}_2) = c(\mathbf{R}_1)\,c(\mathbf{R}_2)$.
The only continuous functions with this property are exponentials:
\begin{equation}
  c(\mathbf{R}) = e^{i\,\mathbf{k}\cdot\mathbf{R}} .
\end{equation}

\paragraph{Step 4: what $\mathbf{k}$ can be.}
\begin{itemize}
  \item If the wave stays bounded over the infinite lattice, then $|c| = 1$
        and $\mathbf{k}$ is \emph{real}. These are the propagating waves,
        and they form the bands.
  \item If $|c| \neq 1$, then $\mathbf{k}$ is \emph{complex}: the wave grows
        or decays from cell to cell (evanescent wave).
\end{itemize}
Keep the second case in mind: inside a band gap no real-$\mathbf{k}$
solutions exist, but complex-$\mathbf{k}$ solutions do. The proof itself does
not restrict $\mathbf{k}$ to real values; restricting to real $\mathbf{k}$ is
our choice, because we want to know which waves propagate.

% ---------------------------------------------------------------------
\section{Consequences}

\begin{itemize}
  \item \textbf{One cell is enough, but only for an infinite lattice.} A finite
        plate has edges; a shift no longer leaves it unchanged, and Step 1
        fails. The band diagram describes the \emph{material}, not the finite
        $1.1 \times 1.0\,\mathrm{m}$ plate.
  \item \textbf{$\mathbf{k}$ only matters modulo $\mathbf{G}$.} The phase
        $e^{i\mathbf{k}\cdot\mathbf{R}}$ does not change when we add a
        reciprocal lattice vector
        $\mathbf{G} = 2\pi\,(m/L_x,\; n/L_y)$, because
        $\mathbf{G}\cdot\mathbf{R}$ is a multiple of $2\pi$. This is the origin
        of the Brillouin zone and of folding.
  \item \textbf{At each $\mathbf{k}$ the problem is finite.} The cell has
        finitely many degrees of freedom, so it has a discrete list of
        frequencies. Sorted, they are bands $1, 2, 3, \dots$ at that
        $\mathbf{k}$.
\end{itemize}

\paragraph{Brillouin zone and symmetry points.}
All distinct wavevectors lie in the first Brillouin zone
$-\pi/L_x \le k_x \le \pi/L_x$, $-\pi/L_y \le k_y \le \pi/L_y$. Reciprocity
($\omega(\mathbf{k}) = \omega(-\mathbf{k})$) and the mirror symmetries of the
cell reduce this to the irreducible Brillouin zone, and band diagrams follow its
edge.

\begin{center}
\begin{tabular}{|l|l|l|l|}
\hline
Point & $\mathbf{k}$ & Phase per cell & Picture \\
\hline
$\Gamma$ & $(0,0)$ & $1,\ 1$ & all cells move identically \\
$X$ & $(\pi/L_x, 0)$ & $-1$ in $x$ & neighbours in $x$ move oppositely \\
$M$ & $(\pi/L_x, \pi/L_y)$ & $-1,\ -1$ & checkerboard \\
$Y$ & $(0, \pi/L_y)$ & $-1$ in $y$ & only needed for rectangular cells \\
\hline
\end{tabular}
\end{center}
Path: $\Gamma$--$X$--$M$--$\Gamma$ for a square cell,
$\Gamma$--$X$--$M$--$Y$--$\Gamma$ for a rectangular cell.

% ---------------------------------------------------------------------
\section{Band folding with our numbers}

For the uniform plate (no cells) each plane wave $e^{i\boldsymbol{\kappa}\cdot\mathbf{x}}$
has one frequency, set by the physical wavenumber $\kappa = |\boldsymbol{\kappa}|$.
Thin-plate (Kirchhoff) estimate for the $5\,\mathrm{mm}$ plate:
\begin{equation}
  f \approx 1.215\,\kappa^2\ \mathrm{Hz}\qquad(\kappa \text{ in rad/m}).
\end{equation}
Described with cells of size $a$, the Bloch wavevector only sees
$\boldsymbol{\kappa}$ modulo $\mathbf{G}$. At a fixed $\mathbf{k}$ the cell
therefore carries all waves $\boldsymbol{\kappa} = \mathbf{k} + \mathbf{G}$,
each with its own frequency. Folding is a relabelling, nothing physical
changes.

\begin{center}
\begin{tabular}{|l|l|l|l|l|}
\hline
Point & $\kappa$ folding there & Kirchhoff & RM (FE) & Multiplicity \\
\hline
$X$ & $\pi/a = 31.4$ & $1199\,\mathrm{Hz}$ & $\approx 1190\,\mathrm{Hz}$ & 2: $(\pm\pi/a, 0)$ \\
$M$ & $\sqrt{2}\,\pi/a = 44.4$ & $2398\,\mathrm{Hz}$ & $\approx 2380\,\mathrm{Hz}$ & 4: $(\pm\pi/a, \pm\pi/a)$ \\
$\Gamma$ & $2\pi/a = 62.8$ & $4797\,\mathrm{Hz}$ & $\approx 4730\,\mathrm{Hz}$ & 4: $(\pm 2\pi/a, 0), (0, \pm 2\pi/a)$ \\
\hline
\end{tabular}
\end{center}
The difference between Kirchhoff and RM grows with frequency
($0.8\,\%$ at $X$, $1.4\,\%$ at $\Gamma$), because shear deformation and
rotary inertia grow with $\kappa h$. This is also why the Kirchhoff
plane-wave model was $1.6\,\%$ off in the stage~4 validation.

% ---------------------------------------------------------------------
\section{Check questions with answers}

\paragraph{Question 1.} \emph{In Step 2, why can we only say that we can
choose solutions that $\mathcal{T}_{\mathbf{R}}$ multiplies by a number,
instead of claiming that every solution is like that?}

\emph{Answer.} Bloch waves are a basis we choose; not every solution is one.
Example: at a generic frequency the plate carries $e^{+i\kappa x}$ and
$e^{-i\kappa x}$, which have \emph{different} phases per cell,
$e^{\pm i\kappa a}$. Their sum $\cos(\kappa x)$ solves the equations at the
same $\omega$, but shifting it by $a$ does not give a multiple of
$\cos(\kappa x)$. It is a valid solution, but not a Bloch wave.

\emph{At $X$ this changes:} $e^{+i\pi x/a}$ and $e^{-i\pi x/a}$ both have the
phase $e^{\pm i\pi} = -1$ per cell. They differ by exactly one $\mathbf{G}$
and are \emph{the same} Bloch $\mathbf{k}$. Every combination of them is
therefore also a Bloch wave at $X$, including $\cos(\pi x/a)$ and
$\sin(\pi x/a)$. This is why the standing waves A and B used in
Chapter~\ref{ch:two_level} are legitimate Bloch modes, and why a resonator can
split them.

\paragraph{Question 2.} \emph{At $\Gamma$ the bare cell has a band at about
$4.73\,\mathrm{kHz}$, although ``$\Gamma$ means infinite wavelength''. What is
that wave physically?}

\emph{Answer.} The Bloch wavevector is $\mathbf{k} = 0$, but the physical
wavenumber is $\kappa = |\mathbf{k} + \mathbf{G}| = 2\pi/a$. Both describe the
same mode because $\mathbf{k}$ only matters modulo $\mathbf{G}$. Physically
it is a wave with wavelength $\lambda_\mathrm{wave} = a = 0.1\,\mathrm{m}$:
exactly one full wiggle per cell, phase per cell $e^{i2\pi} = 1$. All cells
look identical, which is why it sits at $\Gamma$. It appears four times, once
for each direction $\pm x$, $\pm y$.

\paragraph{Question 3.} \emph{What breaks first if Bloch's theorem is applied
to the finite $1.1 \times 1.0\,\mathrm{m}$ plate, and what goes wrong
physically?}

\emph{Answer.} Step 1 already fails: a shifted solution no longer satisfies
the edge conditions, so the operators no longer commute. Physically, waves
reflect at the free edges, travelling waves become standing waves, and the
finite plate has discrete modes. Two consequences (treated later):
\begin{itemize}
  \item the modes of the finite plate roughly \emph{sample} the bands at
        discrete wavenumbers, about $n\pi/L$;
  \item edges can carry localized modes, sometimes even inside a band gap.
\end{itemize}

% =====================================================================

\chapter{Bloch's theorem and band folding}
\label{ch:bloch}

Numbers in this chapter refer to the production configuration of the Bloch
unit cell: steel plate, thickness $h = 5\,\mathrm{mm}$, square unit cell of
size $a = 0.1\,\mathrm{m}$, Reissner--Mindlin (RM) plate theory, bands below
$5\,\mathrm{kHz}$. Plate data: $\rho h \approx 39.3\,\mathrm{kg/m^2}$,
$D \approx 2.3\,\mathrm{kN\,m}$, plate mass per cell
$\rho h a^2 \approx 0.39\,\mathrm{kg}$.

% ---------------------------------------------------------------------
\section{What Bloch's theorem says}

A periodic structure repeats along the lattice vectors
$\mathbf{a}_1 = (L_x, 0)$ and $\mathbf{a}_2 = (0, L_y)$. A lattice vector is
any $\mathbf{R} = m\,\mathbf{a}_1 + n\,\mathbf{a}_2$ with integers $m, n$.
Bloch's theorem states that every free wave can be written in the form
\begin{equation}
  u(\mathbf{x} + \mathbf{R}) = e^{i\,\mathbf{k}\cdot\mathbf{R}}\, u(\mathbf{x}),
  \label{eq:bloch}
\end{equation}
or equivalently $u(\mathbf{x}) = e^{i\mathbf{k}\cdot\mathbf{x}}\,p(\mathbf{x})$
with $p$ periodic in the cell. Moving one cell over does not change the shape
of the wave, it only multiplies it by a phase. Therefore one cell, with the
conditions ``right edge = left edge $\cdot\, e^{i k_x L_x}$'' and
``top edge = bottom edge $\cdot\, e^{i k_y L_y}$'', describes the whole
infinite structure.

% ---------------------------------------------------------------------
\section{Why it holds: the argument in four steps}

\paragraph{Step 1: the structure is unchanged by a lattice shift.}
Define the translation operator $\mathcal{T}_{\mathbf{R}}$ by
$(\mathcal{T}_{\mathbf{R}} u)(\mathbf{x}) = u(\mathbf{x} + \mathbf{R})$.
The equations of motion of the plate have periodic coefficients (the resonator
sits at the same place in every cell). It therefore makes no difference
whether we first shift and then apply the equations, or the other way round:
the two operators \emph{commute}.

\paragraph{Step 2: a shifted solution is again a solution.}
If $u$ is a free wave at frequency $\omega$, then $\mathcal{T}_{\mathbf{R}}u$
is also a free wave at the same $\omega$. Operators that commute share
eigenfunctions. Within the space of solutions at that $\omega$ we can
therefore \emph{choose} solutions that $\mathcal{T}_{\mathbf{R}}$ only
multiplies by a number:
\begin{equation}
  \mathcal{T}_{\mathbf{R}} u = c(\mathbf{R})\, u .
\end{equation}

\paragraph{Step 3: the multiplier must be an exponential.}
Shifting by $\mathbf{R}_1$ and then by $\mathbf{R}_2$ equals shifting by
$\mathbf{R}_1 + \mathbf{R}_2$, so
$c(\mathbf{R}_1 + \mathbf{R}_2) = c(\mathbf{R}_1)\,c(\mathbf{R}_2)$.
The only continuous functions with this property are exponentials:
\begin{equation}
  c(\mathbf{R}) = e^{i\,\mathbf{k}\cdot\mathbf{R}} .
\end{equation}

\paragraph{Step 4: what $\mathbf{k}$ can be.}
\begin{itemize}
  \item If the wave stays bounded over the infinite lattice, then $|c| = 1$
        and $\mathbf{k}$ is \emph{real}. These are the propagating waves,
        and they form the bands.
  \item If $|c| \neq 1$, then $\mathbf{k}$ is \emph{complex}: the wave grows
        or decays from cell to cell (evanescent wave).
\end{itemize}
Keep the second case in mind: inside a band gap no real-$\mathbf{k}$
solutions exist, but complex-$\mathbf{k}$ solutions do. The proof itself does
not restrict $\mathbf{k}$ to real values; restricting to real $\mathbf{k}$ is
our choice, because we want to know which waves propagate.

% ---------------------------------------------------------------------
\section{Consequences}

\begin{itemize}
  \item \textbf{One cell is enough, but only for an infinite lattice.} A finite
        plate has edges; a shift no longer leaves it unchanged, and Step 1
        fails. The band diagram describes the \emph{material}, not the finite
        $1.1 \times 1.0\,\mathrm{m}$ plate.
  \item \textbf{$\mathbf{k}$ only matters modulo $\mathbf{G}$.} The phase
        $e^{i\mathbf{k}\cdot\mathbf{R}}$ does not change when we add a
        reciprocal lattice vector
        $\mathbf{G} = 2\pi\,(m/L_x,\; n/L_y)$, because
        $\mathbf{G}\cdot\mathbf{R}$ is a multiple of $2\pi$. This is the origin
        of the Brillouin zone and of folding.
  \item \textbf{At each $\mathbf{k}$ the problem is finite.} The cell has
        finitely many degrees of freedom, so it has a discrete list of
        frequencies. Sorted, they are bands $1, 2, 3, \dots$ at that
        $\mathbf{k}$.
\end{itemize}

\paragraph{Brillouin zone and symmetry points.}
All distinct wavevectors lie in the first Brillouin zone
$-\pi/L_x \le k_x \le \pi/L_x$, $-\pi/L_y \le k_y \le \pi/L_y$. Reciprocity
($\omega(\mathbf{k}) = \omega(-\mathbf{k})$) and the mirror symmetries of the
cell reduce this to the irreducible Brillouin zone, and band diagrams follow its
edge.

\begin{center}
\begin{tabular}{|l|l|l|l|}
\hline
Point & $\mathbf{k}$ & Phase per cell & Picture \\
\hline
$\Gamma$ & $(0,0)$ & $1,\ 1$ & all cells move identically \\
$X$ & $(\pi/L_x, 0)$ & $-1$ in $x$ & neighbours in $x$ move oppositely \\
$M$ & $(\pi/L_x, \pi/L_y)$ & $-1,\ -1$ & checkerboard \\
$Y$ & $(0, \pi/L_y)$ & $-1$ in $y$ & only needed for rectangular cells \\
\hline
\end{tabular}
\end{center}
Path: $\Gamma$--$X$--$M$--$\Gamma$ for a square cell,
$\Gamma$--$X$--$M$--$Y$--$\Gamma$ for a rectangular cell.

% ---------------------------------------------------------------------
\section{Band folding with our numbers}

For the uniform plate (no cells) each plane wave $e^{i\boldsymbol{\kappa}\cdot\mathbf{x}}$
has one frequency, set by the physical wavenumber $\kappa = |\boldsymbol{\kappa}|$.
Thin-plate (Kirchhoff) estimate for the $5\,\mathrm{mm}$ plate:
\begin{equation}
  f \approx 1.215\,\kappa^2\ \mathrm{Hz}\qquad(\kappa \text{ in rad/m}).
\end{equation}
Described with cells of size $a$, the Bloch wavevector only sees
$\boldsymbol{\kappa}$ modulo $\mathbf{G}$. At a fixed $\mathbf{k}$ the cell
therefore carries all waves $\boldsymbol{\kappa} = \mathbf{k} + \mathbf{G}$,
each with its own frequency. Folding is a relabelling, nothing physical
changes.

\begin{center}
\begin{tabular}{|l|l|l|l|l|}
\hline
Point & $\kappa$ folding there & Kirchhoff & RM (FE) & Multiplicity \\
\hline
$X$ & $\pi/a = 31.4$ & $1199\,\mathrm{Hz}$ & $\approx 1190\,\mathrm{Hz}$ & 2: $(\pm\pi/a, 0)$ \\
$M$ & $\sqrt{2}\,\pi/a = 44.4$ & $2398\,\mathrm{Hz}$ & $\approx 2380\,\mathrm{Hz}$ & 4: $(\pm\pi/a, \pm\pi/a)$ \\
$\Gamma$ & $2\pi/a = 62.8$ & $4797\,\mathrm{Hz}$ & $\approx 4730\,\mathrm{Hz}$ & 4: $(\pm 2\pi/a, 0), (0, \pm 2\pi/a)$ \\
\hline
\end{tabular}
\end{center}
The difference between Kirchhoff and RM grows with frequency
($0.8\,\%$ at $X$, $1.4\,\%$ at $\Gamma$), because shear deformation and
rotary inertia grow with $\kappa h$. This is also why the Kirchhoff
plane-wave model was $1.6\,\%$ off in the stage~4 validation.

% ---------------------------------------------------------------------
\section{Check questions with answers}

\paragraph{Question 1.} \emph{In Step 2, why can we only say that we can
choose solutions that $\mathcal{T}_{\mathbf{R}}$ multiplies by a number,
instead of claiming that every solution is like that?}

\emph{Answer.} Bloch waves are a basis we choose; not every solution is one.
Example: at a generic frequency the plate carries $e^{+i\kappa x}$ and
$e^{-i\kappa x}$, which have \emph{different} phases per cell,
$e^{\pm i\kappa a}$. Their sum $\cos(\kappa x)$ solves the equations at the
same $\omega$, but shifting it by $a$ does not give a multiple of
$\cos(\kappa x)$. It is a valid solution, but not a Bloch wave.

\emph{At $X$ this changes:} $e^{+i\pi x/a}$ and $e^{-i\pi x/a}$ both have the
phase $e^{\pm i\pi} = -1$ per cell. They differ by exactly one $\mathbf{G}$
and are \emph{the same} Bloch $\mathbf{k}$. Every combination of them is
therefore also a Bloch wave at $X$, including $\cos(\pi x/a)$ and
$\sin(\pi x/a)$. This is why the standing waves A and B used in
Chapter~\ref{ch:two_level} are legitimate Bloch modes, and why a resonator can
split them.

\paragraph{Question 2.} \emph{At $\Gamma$ the bare cell has a band at about
$4.73\,\mathrm{kHz}$, although ``$\Gamma$ means infinite wavelength''. What is
that wave physically?}

\emph{Answer.} The Bloch wavevector is $\mathbf{k} = 0$, but the physical
wavenumber is $\kappa = |\mathbf{k} + \mathbf{G}| = 2\pi/a$. Both describe the
same mode because $\mathbf{k}$ only matters modulo $\mathbf{G}$. Physically
it is a wave with wavelength $\lambda_\mathrm{wave} = a = 0.1\,\mathrm{m}$:
exactly one full wiggle per cell, phase per cell $e^{i2\pi} = 1$. All cells
look identical, which is why it sits at $\Gamma$. It appears four times, once
for each direction $\pm x$, $\pm y$.

\paragraph{Question 3.} \emph{What breaks first if Bloch's theorem is applied
to the finite $1.1 \times 1.0\,\mathrm{m}$ plate, and what goes wrong
physically?}

\emph{Answer.} Step 1 already fails: a shifted solution no longer satisfies
the edge conditions, so the operators no longer commute. Physically, waves
reflect at the free edges, travelling waves become standing waves, and the
finite plate has discrete modes. Two consequences (treated later):
\begin{itemize}
  \item the modes of the finite plate roughly \emph{sample} the bands at
        discrete wavenumbers, about $n\pi/L$;
  \item edges can carry localized modes, sometimes even inside a band gap.
\end{itemize}
